\hypertarget{chapter-04-page-089}{}
\subsection{具有偶核的奇异积分}
\label{sec:ch4-even-kernel-singular-integrals}

若奇异积分 \eqref{eq:ch4-singular-integral-definition} 中的 $\Omega$ 是偶函数, 就不能用 Hilbert 变换表示它, 因而旋转方法不适用. 不过, 由 \eqref{eq:ch4-riesz-square-identity}, 应当可以作如下论证:
\[
 Tf=-\sum_{j=1}^nR_j^2(Tf)=-\sum_{j=1}^nR_j(R_jT)f.
\]
算子 $R_jT$ 是奇算子, 因为它是奇算子与偶算子的复合. 如果能证明 $R_jT$ 具有 \eqref{eq:ch4-singular-integral-definition} 的表示, 那么由推论~\ref{cor:ch4-odd-kernel-lp}, $T$ 在 $L^p$ 上有界. 下面把这个论证严格化.

设 $\Omega$ 是 $S^{n-1}$ 上平均值为零的偶函数, 且对某个 $q>1$ 有 $\Omega\in L^q(S^{n-1})$. 对 $\varepsilon>0$, 令
\begin{equation}
\label{eq:ch4-truncated-even-kernel}
 K_\varepsilon(x)=\frac{\Omega(x')}{|x|^n}
 \chi_{\{|x|>\varepsilon\}}.
\end{equation}
注意 $K_\varepsilon\in L^r$, $1<r\le q$. 因而若 $f\in C_c^\infty(\mathbb R^n)$, 取 Fourier 变换可得
\begin{equation}
\label{eq:ch4-riesz-convolution-association}
 R_j(K_\varepsilon*f)=(R_jK_\varepsilon)*f.
\end{equation}

\begin{Lemma}
\label{lem:ch4-even-kernel-riesz-limit}
在上述假设下, 存在奇函数 $\widetilde K_j$, 它是 $-n$ 次齐次的, 并且在每个不含原点的紧集上, 当 $\varepsilon\to0$ 时
\[
 R_jK_\varepsilon(x)\longrightarrow\widetilde K_j(x)
\]
在 $L^\infty$ 范数下成立.
\end{Lemma}

\begin{Proof}
固定 $x\ne0$, 并令 $0<\varepsilon<\nu<|x|/2$. 对几乎处处的此类 $x$, 由推论~\ref{cor:ch4-maximal-truncation-lp},
\[
\begin{aligned}
 R_jK_\varepsilon(x)-R_jK_\nu(x)
 &=c_n\int_{\varepsilon<|y|<\nu}
 \frac{x_j-y_j}{|x-y|^{n+1}}\frac{\Omega(y')}{|y|^n}\,dy\\
 &=c_n\int_{\varepsilon<|y|<\nu}
 \left(\frac{x_j-y_j}{|x-y|^{n+1}}-\frac{x_j}{|x|^{n+1}}\right)
 \frac{\Omega(y')}{|y|^n}\,dy,
\end{aligned}
\]
最后一步用了 $\Omega$ 的平均值为零. 对被积函数应用中值定理, 得
\begin{equation}
\label{eq:ch4-riesz-kernel-cauchy-estimate}
 |R_jK_\varepsilon(x)-R_jK_\nu(x)|
 \le\frac{C}{|x|^{n+1}}
 \int_{\varepsilon<|y|<\nu}\frac{|\Omega(y')|}{|y|^{n-1}}\,dy
 \le\frac{C\|\Omega\|_1}{|x|^{n+1}}\nu.
\end{equation}
因此, 对每个 $a>0$, $\{R_jK_\varepsilon\}$ 在 $\{|x|>a\}$ 上按 $L^\infty$ 范数是 Cauchy 列. 将其几乎处处极限记为 $K_j^*$. 因为 $R_jK_\varepsilon$ 是奇函数, 可以把 $K_j^*$ 取成奇函数.

由伸缩关系可知 $K_j^*(\lambda x)=\lambda^{-n}K_j^*(x)$ 对几乎处处的 $(x,\lambda)$ 成立. 通过在某个以原点为中心的球面上选取一个零测集外的代表, 再沿射线按 $-n$ 次齐次方式定义, 得到处处齐次的可测奇函数 $\widetilde K_j$, 且 $\widetilde K_j=K_j^*$ 几乎处处. 估计 \eqref{eq:ch4-riesz-kernel-cauchy-estimate} 随即给出所断言的局部一致收敛.
\end{Proof}

\hypertarget{chapter-04-page-090}{}
\begin{Lemma}
\label{lem:ch4-even-kernel-l1-estimates}
引理~\ref{lem:ch4-even-kernel-riesz-limit} 中的核 $\widetilde K_j$ 满足
\[
 \int_{S^{n-1}}|\widetilde K_j(u)|\,d\sigma(u)
 \le C_q\|\Omega\|_q.
\]
此外, 若 $\widetilde K_{j,\varepsilon}(x)=\widetilde K_j(x)
\chi_{\{|x|>\varepsilon\}}$, 且
$\Delta_\varepsilon=R_jK_\varepsilon-\widetilde K_{j,\varepsilon}$, 则 $\Delta_\varepsilon\in L^1(\mathbb R^n)$, 并且
$\|\Delta_\varepsilon\|_1\le C_q\|\Omega\|_q$.
\end{Lemma}

\begin{Proof}
由 $\widetilde K_j$ 的齐次性,
\[
\begin{aligned}
 \int_{S^{n-1}}|\widetilde K_j(u)|\,d\sigma(u)
 &=\frac1{\log2}\int_{1<|x|<2}|\widetilde K_j(x)|\,dx\\
 &\le\frac1{\log2}\int_{1<|x|<2}
 |\widetilde K_j(x)-R_jK_{1/2}(x)|\,dx
 +\frac1{\log2}\int_{1<|x|<2}|R_jK_{1/2}(x)|\,dx.
\end{aligned}
\]
令 \eqref{eq:ch4-riesz-kernel-cauchy-estimate} 中 $\nu=1/2$ 并令 $\varepsilon\to0$, 得
\begin{equation}
\label{eq:ch4-riesz-kernel-pointwise-estimate}
 |\widetilde K_j(x)-R_jK_{1/2}(x)|
 \le\frac{C\|\Omega\|_1}{|x|^{n+1}},
 \qquad |x|>1.
\end{equation}
第一项因而由 $C\|\Omega\|_1$ 控制. 第二项由 Riesz 变换的 $L^q$ 有界性控制为 $C\|\Omega\|_q$. 这就证明了第一个断言.

由伸缩, $\Delta_\varepsilon(x)=\varepsilon^{-n}\Delta_1(\varepsilon^{-1}x)$, 因而只需证明 $\|\Delta_1\|_1<\infty$. 将积分分成 $|x|<2$, $1<|x|<2$ 与 $|x|>2$ 三部分. 前两部分分别由 Riesz 变换的 $L^q$ 有界性和刚证明的球面估计控制. 对第三部分重复 \eqref{eq:ch4-riesz-kernel-pointwise-estimate} 的论证, 得 $|\Delta_1(x)|\le C\|\Omega\|_1|x|^{-n-1}$, 从而完成证明.
\end{Proof}

\hypertarget{chapter-04-page-091}{}
\begin{Theorem}
\label{thm:ch4-rough-kernel-lp}
设 $\Omega$ 是 $S^{n-1}$ 上平均值为零的函数, 它的奇部属于 $L^1(S^{n-1})$, 偶部属于某个 $L^q(S^{n-1})$, $q>1$. 则 \eqref{eq:ch4-singular-integral-definition} 中的奇异积分 $T$ 在 $L^p(\mathbb R^n)$, $1<p<\infty$, 上有界.
\end{Theorem}

\begin{Proof}
由推论~\ref{cor:ch4-odd-kernel-lp}, 只需考虑 $\Omega$ 为偶函数的情形. 与 Hilbert 变换的论证相同, 只需对 $f\in C_c^\infty(\mathbb R^n)$ 建立 $L^p$ 不等式. 对这样的 $f$, $Tf=\lim_{\varepsilon\to0}K_\varepsilon*f$. 由 \eqref{eq:ch4-riesz-square-identity} 与 \eqref{eq:ch4-riesz-convolution-association},
\[
 K_\varepsilon*f=-\sum_{j=1}^nR_j((R_jK_\varepsilon)*f).
\]
在引理~\ref{lem:ch4-even-kernel-l1-estimates} 的记号下,
\[
 (R_jK_\varepsilon)*f=\widetilde K_{j,\varepsilon}*f+\Delta_\varepsilon*f.
\]
引理~\ref{lem:ch4-even-kernel-riesz-limit} 给出奇的 $-n$ 次齐次核 $\widetilde K_j$, 因而由推论~\ref{cor:ch4-maximal-truncation-lp} 和引理~\ref{lem:ch4-even-kernel-l1-estimates},
\[
 \|\widetilde K_{j,\varepsilon}*f\|_p
 \le C\int_{S^{n-1}}|\widetilde K_j(u)|\,d\sigma(u)\|f\|_p
 \le C\|\Omega\|_q\|f\|_p,
\]
且
\[
 \|\Delta_\varepsilon*f\|_p
 \le\|\Delta_\varepsilon\|_1\|f\|_p
 \le C\|\Omega\|_q\|f\|_p.
\]
结合这些估计及 $R_j$ 在 $L^p$ 上的有界性, 得 $\|K_\varepsilon*f\|_p\le C\|\Omega\|_q\|f\|_p$. 右端与 $\varepsilon$ 无关, Fatou 引理遂给出同样的 $Tf$ 估计.
\end{Proof}

\hypertarget{chapter-04-page-092}{}
定理~\ref{thm:ch4-rough-kernel-lp} 的另一个证明将在\MBUnresolvedReference{chapter}{第 8 章}给出, 见\MBUnresolvedReference{corollary}{推论 8.21}.
